\section{A cofinal square multiple for the beta modulus}
\label{sec:odd-square-beta}

This section independently verifies a further change to the factorial-free
certificate.  It concerns the beta modulus parameter $y$, not the radix $q$
or the generalized binomial CRT modulus $B$.

Let $t\ge16$ and let
\[
 F=\bigl(2(t+1)\bigr)^t.
\]
The old front end imposed $y\ge F+3$ by a ternary nonnegative-value
certificate.  Replace that certificate by two quadratic equations in two
unrestricted integer coordinates $u,w$:
\begin{equation}\label{eq:odd-square-beta-rows}
 w-4u^2-4u-1=0,\qquad y-Fw=0.
\end{equation}
Thus $w=(2u+1)^2\ge1$ and $y\ge F>0$.  No factorial is certified by
these equations.  Factorials used below occur only in a finite existence
proof.

\subsection{The initial sign and size recovery remains acyclic}

Use the selected factorial-free front end: $q$ is an odd square;
$H=2+\sum_j a_j$; all nine actual representatives $a_j$ are
nonnegative by their requests at base $q+1\ge2$; and the index row
$C_0+F^2+1=B\iota$ implies $B\ne0$.  Therefore $|BH|\ge2$.
The signed power certificate for $q=(BH)^x$ excludes $x<0$, so $x\ge0$.
This reasoning does not require $B>0$ or $q>1$ in advance.

Retain the two endpoint equations, with unrestricted integer quotients:
\begin{equation}\label{eq:odd-square-beta-endpoints}
 x=F+1+(y+1)u_0,\qquad
 x=1+\bigl(1+(t+1)y\bigr)u_t.
\end{equation}
The second equation excludes $x=0$: otherwise the integer
$1+(t+1)y>1$ would divide $1$.  Hence $x\ge1$.
If $u_0\le-1$, the first equation would give
$x\le F-y\le0$, a contradiction.  It follows that
\[
 u_0\ge0,\qquad x\ge F+1>1.
\]
The second endpoint then gives $u_t\ge1$ and
\begin{equation}\label{eq:odd-square-beta-x-bound}
 x\ge2+(t+1)y\ge2+(t+1)F.
\end{equation}

These are sufficient for every earlier bootstrap bound.  In particular,
for $F\ge7$ and $t\le F$,
\[
 F^2+t\le F(F+1)\le2^{F-1}.
\]
The last inequality holds at $F=7$ and is preserved upon increasing $F$,
since $(F+2)/F<2$.  For $0\le j\le t$ it follows that
\[
 \binom{F^2+t}{j}\le(F^2+t)^t\le2^{t(F-1)}
 \le q/2^{F+t+2}<q/64.
\]
Here $q=|BH|^x\ge2^x$ and
$x-t(F-1)\ge F+t+2>6$.  In particular the largest relevant coefficient
divided by $q-1$ is less than $1/32$.  The geometric-series estimate
therefore places the ordinary low remainder strictly below $q^t/32$.
Thus the first centered binomial digit can be decoded before any
interpretation of the local beta rows, with enough margin for the common
remainder budget.

The local arithmetic majorant is unchanged.  If
$R_0=2+t+x+y+F$, then \eqref{eq:odd-square-beta-x-bound} gives
$R_0<2x$, because
$t(y-1)>F$.  Also $x\ge F+1>2^{32}$.  The inherited exponent
$E=2^{23}$ therefore still gives
$L<R_0^E<(2x)^E<2^x$.  The lower indices $t,t+1,y+1,x$ remain at most
$K=x$.  Finally, $F>2^{22}$ gives
$2^{22}(t+1)<x$, retaining the derived-label digit bound.  None of
these arguments presupposes that $y>F$.

\subsection{The exceptional boundary $y=F$ is impossible at a zero}

After decoding the generalized CRT certificate, obtain the same actual
local source rows as before.  Denote their accumulators by $S_i,S'_i$.
Their range conditions give
\[
 0\le S_i,S'_i\le y,
\]
and the beta rows use the consecutive moduli
$D_i=1+(i+1)y$.  Because $y<D_i$ for every $i\ge0$, ordinary
remainder uniqueness implies $S'_i=S_{i+1}$ for $i<t-1$.

Suppose $y=F$.  The first endpoint says that $x$ is divisible by
$y+1$, hence $S_0=0$.  Every selected source rule has a flow identity
\[
 S'_i(1+v_iF)=S_i(1+p_iF)(1+q_iF).
\]
The label weights are positive integers.  Consequently $S_i=0$ implies
$S'_i=0$, and the overlaps propagate zero through every row.  However,
the last endpoint says $x\equiv1\pmod{1+(t+1)y}$, so the final
bounded remainder is $S'_{t-1}=1$.  This is a contradiction.

Therefore every zero has $y>F$.  Indeed the odd-square parametrization
gives the stronger $y\ge9F$.  In particular $F+1\le y$, and the first
endpoint now gives $S_0=F+1$ by ordinary remainder uniqueness.  The
usual multiplication and cancellation of positive accumulators, followed
by polynomial multiplicity separation, applies without any further
change.  This boundary exclusion is essential: the inequality $y\ge F$
alone does not make $F+1$ a canonical initial remainder.

\subsection{Completeness and freedom to choose the parity of $x$}

Start from a finite proof tree.  First choose the permitted horizon $t$
large enough for its step count, labels, parameters and padding rows.
Then fix $F$ and all finitely many pending-multiset products.

Write $t!=2^a m$ with $m$ odd.  Since
\[
 a=v_2(t!)<t,\qquad v_2(F)\ge t,
\]
the integer $F$ supplies the entire power of two in $t!$.  Choose an
arbitrarily large positive odd multiple $v$ of $m$, and put
\[
 u=(v-1)/2,\qquad w=v^2,\qquad y=Fv^2.
\]
Then $t!\mid y$, $y$ is even, and $y$ can exceed every pending product
and $F+1$.  Both rows \eqref{eq:odd-square-beta-rows} hold.

The moduli $1+(i+1)y$, $0\le i\le t$, are pairwise coprime.  A common
prime divisor of two of them would divide their index difference, hence
would be at most $t$ and would divide $t!\mid y$, which is impossible.
CRT supplies the prescribed accumulator residues.  Since every modulus
is odd, the common representative $x$ may be chosen even, positive, and
arbitrarily large.  Its quotients are nonnegative and its row quotients
are strictly less than $x$.

The remaining front variables, generalized binomial CRT representatives,
and private Pell witnesses are then chosen in their established order.
In particular the parity choices making $H$ odd and the semantic anchor
even remain independent of this beta-CRT step.  Thus $q=(BH)^x$ is an
odd square, as required.  This proves completeness of the replacement.

\paragraph{Exact cost boundary.}
Before any subsequent projection, this replaces three ternary-guard
witnesses and one quadratic row by two witnesses and two quadratic rows.
It therefore saves one unrestricted integer coordinate at degree four,
while adding one quadratic residual.  Counts after algebraic contraction
must be read from the emitted final equation, not inferred from this
local accounting alone.
